% Problem record op_d2813fe3fcdf09ad. This TeX file is derived from
% database/problems_json/op_d2813fe3fcdf09ad.json by scripts/sync-tex.mjs; edit the JSON record
% and rerun the script rather than editing this file.
\section{Private capacity of the qubit depolarizing channel}
\label{sec:op_d2813fe3fcdf09ad}

\paragraph{Problem.}
What is the private classical capacity $P(\Lambda_p)$ of the qubit depolarizing
channel, defined for a mixing parameter $p\in[0,1]$ by
\begin{equation}
  \Lambda_p(\rho):=(1-p)\,\rho+p\,\mathrm{Tr}[\rho]\,\frac{I}{2}
  =\bigl(1-\tfrac{3p}{4}\bigr)\rho
  +\tfrac{p}{4}\bigl(X\rho X+Y\rho Y+Z\rho Z\bigr),
  \label{eq:pd-depolarizing-channel}
\end{equation}
where $\rho$ ranges over qubit states, $I$ is the qubit identity, and
$X,Y,Z$ are the Pauli matrices?  The second expression in
Eq.~\eqref{eq:pd-depolarizing-channel} exhibits the Pauli error probabilities
$\bigl(p_I,p_X,p_Y,p_Z\bigr)=\bigl(1-\tfrac{3p}{4},\tfrac{p}{4},\tfrac{p}{4},
\tfrac{p}{4}\bigr)$, with total Pauli error probability $\tfrac{3p}{4}$.
Fix a Stinespring isometry $V:A\to BE$ of $\Lambda_p$ and a finite ensemble
$\{q_x,\rho_x\}_x$ of qubit states, and set
$\omega^{XBE}:=\sum_x q_x\,|x\rangle\!\langle x|^X\otimes V\rho_xV^\dagger$.
The one-shot private information and the private classical capacity of the
channel are
\begin{equation}
  P^{(1)}(\Lambda_p):=\max_{\{q_x,\rho_x\}}
  \bigl[I(X;B)_\omega-I(X;E)_\omega\bigr],
  \qquad
  P(\Lambda_p):=\sup_{n\ge1}\frac{1}{n}\,
  P^{(1)}\bigl(\Lambda_p^{\otimes n}\bigr),
  \label{eq:pd-private-capacity}
\end{equation}
where $I(\cdot\,;\cdot)_\omega$ is the quantum mutual information and all
logarithms are base 2 \sourcecite{ref:pd-devetak}{Dev05},
\sourcecite{ref:pd-cai-winter-yeung}{CWY04}.  Determine the value of
$P(\Lambda_p)$ in Eq.~\eqref{eq:pd-private-capacity} as an explicit function
of $p$ on $[0,1]$.  For context, every channel satisfies $P\ge\mathcal Q$,
where $\mathcal Q$ denotes the quantum capacity, the supremum of achievable
coherent-information rates, because the coherent information of any input
state equals the private information of the ensemble formed from a
purification of that state \sourcecite{ref:pd-devetak}{Dev05}; whether
$P(\Lambda_p)$ and $\mathcal Q(\Lambda_p)$ differ anywhere on
$0<p<\tfrac{1}{3}$, i.e. whether shielding the message from the environment
with classical side information strictly helps on this channel, is part of
what is open.

\subsection*{Status}

Unsolved

\subsection*{Source}

The question is stated explicitly by Leditzky, Leung, and Smith, who note that
it \emph{has been an open problem for more than 20 years to determine the
capacities of some of these low-noise channels such as the depolarizing
channel} \sourcecite{ref:pd-leditzky-leung-smith}{LLS18}.  The regularized
private information defining the capacity is due to Devetak and, in the
wiretap formulation, to Cai, Winter, and Yeung
\sourcecite{ref:pd-devetak}{Dev05},
\sourcecite{ref:pd-cai-winter-yeung}{CWY04}.

\subsection*{Progress}

\begin{itemize}
  \item The channel in Eq.~\eqref{eq:pd-depolarizing-channel} is antidegradable
  precisely for $p\ge\tfrac{1}{3}$, and an antidegradable channel has
  vanishing private capacity because the environment can reconstruct
  everything the receiver obtains, so $P(\Lambda_p)=0$ throughout
  $p\ge\tfrac{1}{3}$.  Kondra et al. prove for the quantum capacity an
  all-code exponential strong converse that, in the convention of
  Eq.~\eqref{eq:pd-depolarizing-channel}, reads
  $\mathcal Q(\Lambda_p)\le\max\{1-3p,0\}$ and vanishes at $p=\tfrac{1}{3}$;
  their result concerns $\mathcal Q$ and leaves $P(\Lambda_p)$ undetermined
  for every $p<\tfrac{1}{3}$ \sourcecite{ref:pd-kondra-et-al}{KBK+26}.

  \item The capacity is bounded below by the quantum capacity and hence by the
  hashing rate: for the Pauli error vector of
  Eq.~\eqref{eq:pd-depolarizing-channel},
  \begin{equation}
    P(\Lambda_p)\ \ge\ \mathcal Q(\Lambda_p)\ \ge
    \max\Bigl\{0,\;1-H_2\bigl(\tfrac{3p}{4}\bigr)
    -\tfrac{3p}{4}\log_2 3\Bigr\},
    \qquad
    H_2(q):=-q\log_2 q-(1-q)\log_2(1-q),
    \label{eq:pd-hashing-bound}
  \end{equation}
  where $H_2$ is the binary entropy and the second inequality is the coherent
  information at the maximally mixed input, itself a private-information rate
  \sourcecite{ref:pd-devetak}{Dev05}.  The bracket in
  Eq.~\eqref{eq:pd-hashing-bound} is positive below the vanishing point
  $\tfrac{3p}{4}\approx0.18929$ and negative above it.

  \item Positivity beyond Eq.~\eqref{eq:pd-hashing-bound} is certified through the
  quantum capacity: Krohn-Grimberghe exhibits a rank-two,
  permutation-symmetric input across $45$ channel uses whose coherent
  information is provably positive at per-Pauli error $16239/250000$, so
  $\mathcal Q(\Lambda_p)>0$, and therefore $P(\Lambda_p)>0$, whenever
  $\tfrac{p}{4}\le\tfrac{16239}{250000}$, that is $p\le0.259824$
  \sourcecite{ref:pd-krohn-grimberghe}{KG26}.  Whether $P(\Lambda_p)>0$
  throughout $(0.259824,\tfrac{1}{3})$ remains open.

  \item For low noise both capacities are known to leading order: Leditzky, Leung,
  and Smith show that for channels close to the identity, including
  $\Lambda_p$ at small $p$, the private and quantum capacities each equal the
  single-letter coherent information $\max_\rho I_c(\rho,\Lambda_p)$ to
  leading order in the distance to the identity channel, so shielding part of
  the data with the auxiliary classical register gains nothing to that order;
  the exact values are left open
  \sourcecite{ref:pd-leditzky-leung-smith}{LLS18}.

  \item The regularization in Eq.~\eqref{eq:pd-private-capacity} cannot be dropped
  in general: the private capacity is not additive
  \sourcecite{ref:pd-li-winter-zou-guo}{LWZG09}, the private information is
  superadditive for every number of channel uses
  \sourcecite{ref:pd-elkouss-strelchuk}{ES15}, and for depolarizing channels
  of growing dimension the superadditivity of the quantum capacity provably
  weakens \sourcecite{ref:pd-etxezarreta-martinez}{EMC+23}.  No blocklength
  $n(p)$ at which the regularization stabilizes is known for this channel.

  \item Since the channel is degradable only at $p=0$
  \sourcecite{ref:pd-cubitt-ruskai-smith}{CRS08}, no single-letter
  degradability formula applies, and the strongest upper bounds come from
  flagged extensions.  Kianvash, Fanizza, and Giovannetti construct
  degradable flag extensions that give the state-of-the-art upper bounds on
  the quantum and private capacities of the depolarizing channel
  \sourcecite{ref:pd-kianvash-fanizza-giovannetti}{KFG22}; Poshtvan and
  Karimipour upper-bound both capacities of the SO(2)-covariant Pauli family
  containing the channel of Eq.~\eqref{eq:pd-depolarizing-channel} across its
  whole parameter space \sourcecite{ref:pd-poshtvan-karimipour}{PK22}; and
  Nourozi refines the flagged-extension technique into single-letter upper
  bounds on both capacities with numerical evaluation, describing the
  depolarizing channel as one whose capacity genuinely requires multi-letter
  optimization \sourcecite{ref:pd-nourozi}{Nour26}.  No known upper bound
  meets a lower bound at any fixed $p\in(0,\tfrac{1}{3})$.

  \item An additive upper bound of a different kind comes from squashed
  entanglement: $P(\Lambda_p)\le P_2(\Lambda_p)\le E_{\mathrm{sq}}(\Lambda_p)$,
  where $P_2$ is the two-way-assisted private capacity and
  $E_{\mathrm{sq}}$ the squashed entanglement of the channel, which is
  additive over tensor products
  \sourcecite{ref:pd-takeoka-guha-wilde}{TGW14}.  The exact value of
  $E_{\mathrm{sq}}(\Lambda_p)$ is itself unknown; it is archived as a
  separate problem (see Comment).
\end{itemize}

\subsection*{References}

\begin{enumerate}[label={},leftmargin=4.8em,labelsep=0.5em]
  \footnotesize
  \item[\textup{[Dev05]}]\label{ref:pd-devetak}
  I. Devetak, ``The Private Classical Capacity and Quantum Capacity of a
  Quantum Channel,'' \emph{IEEE Transactions on Information Theory}
  \textbf{51}, 44--55 (2005).
  \href{https://doi.org/10.1109/TIT.2004.839515}{doi:10.1109/TIT.2004.839515};
  \href{https://arxiv.org/abs/quant-ph/0304127}{arXiv:quant-ph/0304127}.

  \item[\textup{[CWY04]}]\label{ref:pd-cai-winter-yeung}
  N. Cai, A. Winter, and R. W. Yeung, ``Quantum privacy and quantum wiretap
  channels,'' \emph{Problems of Information Transmission} \textbf{40},
  318--336 (2004).
  \href{https://doi.org/10.1007/s11122-005-0002-x}{doi:10.1007/s11122-005-0002-x}.

  \item[\textup{[CRS08]}]\label{ref:pd-cubitt-ruskai-smith}
  T. S. Cubitt, M.-B. Ruskai, and G. Smith, ``The structure of degradable
  quantum channels,'' \emph{Journal of Mathematical Physics} \textbf{49},
  102104 (2008).
  \href{https://doi.org/10.1063/1.2953685}{doi:10.1063/1.2953685};
  \href{https://arxiv.org/abs/0802.1360}{arXiv:0802.1360}.

  \item[\textup{[LLS18]}]\label{ref:pd-leditzky-leung-smith}
  F. Leditzky, D. Leung, and G. Smith, ``Quantum and private capacities of
  low-noise channels,'' \emph{Physical Review Letters} \textbf{120},
  160503 (2018).
  \href{https://doi.org/10.1103/PhysRevLett.120.160503}{doi:10.1103/PhysRevLett.120.160503};
  \href{https://arxiv.org/abs/1705.04335}{arXiv:1705.04335}.

  \item[\textup{[KFG22]}]\label{ref:pd-kianvash-fanizza-giovannetti}
  F. Kianvash, M. Fanizza, and V. Giovannetti, ``Bounding the quantum
  capacity with flagged extensions,'' \emph{Quantum} \textbf{6}, 647
  (2022).
  \href{https://doi.org/10.22331/q-2022-02-09-647}{doi:10.22331/q-2022-02-09-647};
  \href{https://arxiv.org/abs/2008.02461}{arXiv:2008.02461}.

  \item[\textup{[TGW14]}]\label{ref:pd-takeoka-guha-wilde}
  M. Takeoka, S. Guha, and M. M. Wilde, ``The squashed entanglement of a
  quantum channel,'' \emph{IEEE Transactions on Information Theory}
  \textbf{60}(8), 4987--4998 (2014).
  \href{https://doi.org/10.1109/TIT.2014.2330313}{doi:10.1109/TIT.2014.2330313};
  \href{https://arxiv.org/abs/1310.0129}{arXiv:1310.0129}.

  \item[\textup{[PK22]}]\label{ref:pd-poshtvan-karimipour}
  A. Poshtvan and V. Karimipour, ``Capacities of the covariant Pauli
  channel,'' \emph{Physical Review A} \textbf{106}, 062408 (2022).
  \href{https://doi.org/10.1103/PhysRevA.106.062408}{doi:10.1103/PhysRevA.106.062408};
  \href{https://arxiv.org/abs/2206.06106}{arXiv:2206.06106}.

  \item[\textup{[LWZG09]}]\label{ref:pd-li-winter-zou-guo}
  K. Li, A. Winter, X. Zou, and G.-C. Guo, ``The private capacity of
  quantum channels is not additive,'' \emph{Physical Review Letters}
  \textbf{103}, 120501 (2009).
  \href{https://doi.org/10.1103/PhysRevLett.103.120501}{doi:10.1103/PhysRevLett.103.120501};
  \href{https://arxiv.org/abs/0903.4308}{arXiv:0903.4308}.

  \item[\textup{[ES15]}]\label{ref:pd-elkouss-strelchuk}
  D. Elkouss and S. Strelchuk, ``Superadditivity of private information
  for any number of uses of the channel,'' \emph{Physical Review Letters}
  \textbf{115}, 040501 (2015).
  \href{https://doi.org/10.1103/PhysRevLett.115.040501}{doi:10.1103/PhysRevLett.115.040501};
  \href{https://arxiv.org/abs/1502.05326}{arXiv:1502.05326}.

  \item[\textup{[EMC+23]}]\label{ref:pd-etxezarreta-martinez}
  J. Etxezarreta Martinez, A. de Marti i Olius, and P. M. Crespo,
  ``The superadditivity effects of quantum capacity decrease with
  the dimension for qudit depolarizing channels,'' \emph{Physical Review A}
  \textbf{108}, 032602 (2023).
  \href{https://doi.org/10.1103/PhysRevA.108.032602}{doi:10.1103/PhysRevA.108.032602};
  \href{https://arxiv.org/abs/2301.10132}{arXiv:2301.10132}.

  \item[\textup{[Nour26]}]\label{ref:pd-nourozi}
  V. Nourozi, ``Flagged Extensions and Numerical Simulations for Quantum
  Channel Capacity: Bridging Theory and Computation,'' arXiv preprint
  (2026).
  \href{https://arxiv.org/abs/2506.03429}{arXiv:2506.03429}.

  \item[\textup{[KBK+26]}]\label{ref:pd-kondra-et-al}
  T. V. Kondra, R. Brinster, H. Kampermann, D. Bru{\ss}, and N. Wyderka,
  ``Sharp Quantum Capacity Thresholds: Exponential Strong Converses for
  Degradable and Antidegradable Channels,'' arXiv preprint (2026).
  \href{https://arxiv.org/abs/2608.01308}{arXiv:2608.01308}.

  \item[\textup{[KG26]}]\label{ref:pd-krohn-grimberghe}
  A. Krohn-Grimberghe, ``A Certified Lower Bound on the Quantum-Capacity
  Threshold of the Depolarizing Channel,'' arXiv preprint (2026).
  \href{https://arxiv.org/abs/2608.15870}{arXiv:2608.15870}.
\end{enumerate}

\subsection*{Comment}

Beyond the zero region $p\ge\tfrac{1}{3}$, the certified positivity region
$p\le0.259824$, and the leading-order behavior as $p\to0$, nothing exact is
known about $P(\Lambda_p)$: no upper bound meets a lower bound at any fixed
$p\in(0,\tfrac{1}{3})$; positivity on $(0.259824,\tfrac{1}{3})$ is open; it
is open whether the regularization in Eq.~\eqref{eq:pd-private-capacity}
stabilizes at some finite blocklength $n(p)$, the negative answer to the
universal-truncation question for channel-independent integers (solved record
\texttt{01M1Q787QR8FTR00QF4PMHHWPE}) not excluding finite stabilization for
this particular channel; and it is open whether
$P(\Lambda_p)>\mathcal Q(\Lambda_p)$ somewhere.  The sibling record
\texttt{01M20868F4GEPX6FBJZF8T0GRJ} asks for the squashed entanglement
$E_{\mathrm{sq}}(\Lambda_p)$ of the same channel: it upper-bounds the
assisted capacity $P_2(\Lambda_p)\ge P(\Lambda_p)$, but determining it would
neither determine $P(\Lambda_p)$ nor be determined by it.  The record
\texttt{01M1HME7809M71BG24CSMYKA8A} asks for the quantum capacity of the
same channel family; $P\ge\mathcal Q$ is the only known exact
relation between the two values, equality being unproven for every
$p>0$.

\subsection*{Field}

Quantum Communication; Quantum Cryptography

\subsection*{Topic}

Private capacity; Additivity and regularization

\subsection*{ID}

\texttt{op\_d2813fe3fcdf09ad}
